\scanpage{38}
Je me \unsure{rapporte|repporte|reporte|reviens} au diagramme (23)

\diagram{\begin{tikzcd}[ampersand replacement=\&,column sep=small,row sep=small]
\& \& \& 1\arrow[dl] \& \&\\
\& \& \mathbb F_2^{\underline\omega}\arrow[dl]\arrow[drr,equal] \& \& \& 1\arrow[dl]\\
\& \mathrm{Gl}_{\Lambda}(M)\arrow[dl]\arrow[dr] \& \& \& \mathbb F_2^{\underline\omega}\simeq V(S_Q)\arrow[dll] \&\\
\mathfrak S_I\arrow[dr] \& \text{cart} \& \mathcal D_Q\arrow[dl,"{\chi_C}"'] \& \& \&\\
\& \{\pm1\}\arrow[dl] \& \& \& \&\\
\unsure{1|\mathfrak z} \& \& \& \& \&
\end{tikzcd}}

qui exprime $\mathrm{Gl}_\Lambda(M)$ comme produit fibré, la donnée de $M$ équivalant à celle du triple $(I,Q,k)$, où $I$ est un ens. à quatre éléments, $Q$ un carré \inserted{combinatoire}, et $k$ un isom

\[
k:\quad\underbrace{C_Q}_{\text{codiagonales}}\quad\overset{\sim}{\rule[.2ex]{1.8em}{.4pt}}\quad\underline\omega_I
\]

En termes de $Q$, on a

\[
 C_\Delta\simeq F(\Delta)/M_0\simeq P_{M,\mathbb F_2^{\underline\omega}}/\mathbb F_2=S_Q/\mathbb F_2\simeq D_Q
\]

ens des diagonales de $Q$ -- donc la donnée de l'octaèdre $\Delta$ (au lieu de $M\longleftrightarrow(I,Q,k)$) équivaut à celle de $I$, et des deux ens. à deux éléments $C$ ($\simeq\omega_I$) et $D$ correspondant aux diagonales et codiagonales du carré $Q$, \deleted{\unsure{ayant|dont|alors que}} la donnée équivaut à celle d'un torseur sous $\deleted{\unsure{\mathcal D_Q}}\mathbb D_4/\mathfrak z(\mathbb D_4)\simeq\deleted{\unsure{\mathbb F_2|\mathbb D_2}}\mathbb F_2\times\mathbb F_2$, qu'il faut donc relever en un torseur sous $\mathbb D_4$. Ou encore, décrire la donnée d'un carré combinatoire, dont on a déjà les deux ensembles \inserted{$C$ et $D$} des codiagonales et diagonales\ldots\ $]\!]$
% archive sheet number: 37
